\documentclass{article}

\usepackage{csquotes}
\usepackage[backend=biber,natbib=true,style=numeric,sorting=none]{biblatex}

\usepackage{mdframed}

\usepackage{geometry}
 \geometry{
 a4paper,
 total={170mm,257mm},
 left=20mm,
 top=20mm,
 }
\usepackage[english]{babel}
\usepackage{amsthm}
\theoremstyle{definition}
\newtheorem{definition}{Definition}[section]

\usepackage[acronym]{glossaries}
\usepackage{mathtools}% Loads amsmath
\numberwithin{equation}{subsection}
\usepackage{optidef}
\usepackage{cancel}

\addbibresource{LitRevBib.bib}
\usepackage{hyperref} 

\usepackage{tabularx}

\makeglossaries

\newglossaryentry{entryOne}
{
        name=Glossary Entry,
        description={Glossary entries are used to provide definitions for words in your document}
}    

\newcommand{\de}{\mathrm{d}}
\newcommand{\De}{\mathrm{D}}
\newcommand{\nl}{\par\noindent}
\newcommand{\del}{\partial}

\title{Literature Review}
\author{Kazuhide Ogawa Javes 1462313}
\date{September 2026}

%---------------------------------------------------------------------------------------
% BEGIN DOCUMENT HERE
%---------------------------------------------------------------------------------------

\begin{document}
\maketitle

\section{Introduction}
After discussion with A/Prof Douglas R. Brumley, My literature review will be around ciliary and flagellar flow.

\section{Ciliary and flagellar flow}
\subsection{Definitions}
\begin{mdframed}
\begin{definition}[Cilium (pl. Cilia)]
\textit{Cilia} are short, hair-like protrusions that grow from the surface as an extension of the membrane. \cite{mizunoStructuralStudiesCiliary2012}
\end{definition}
\end{mdframed}

\begin{mdframed}
\begin{definition}[Flagellum (pl. Flagellum)]
\textit{Flagella} are whip-like appendages that are found in many cells and organisms such as bacteria, sperm cells and algae.
\cite{brumleyFlagellarSynchronizationDirect2014}
\end{definition}
\end{mdframed}
\textit{Beating} - the periodic, circular motion of cilia and flagella - is a fundamental behaviour which has been extensively studied since at least (insert year). They serve important functions in motility, sensory reception, signalling, and mass transport. Studies of ciliary and flagellar flows therefore have significant biological implications. \cite{mizunoStructuralStudiesCiliary2012}

\subsection{Corals}
Cilia-induced vortical flows are critical for $\text{O}_2$ regulation and metabolite exchange across coral-water interfaces. \nl
\textbf{Pacherres et al.} \cite{pacherresAcuteTemperatureEffects}
\begin{itemize}
        \item Methodology: "high-speed imaging with a microscope chamber to determine cilia beating frequencies (CBFs) and particle image velocimetry (PIV) in combination with chemical imaging (SensPIV) for mapping flow and O2 concentration fields of corals exposed to increasing seawater temperature."
        
        \item Found that ciliary beating is important for thermal tolerance for corals. By beating with higher frequency, the corals increase $\text{O}_2$ intake, which alleviates hypoxia observed in temperatures $\sim 35 \text{C}^{\circ}$ and above. Above $\sim 37 \text{C}^{\circ}$, coral morality accelerated, indicating the existence of a physiological tipping-point for reef-building corals.
\end{itemize}
\nl
A shortcoming of this study and its conclusions is that it is unclear what specific temperature-dependent dynamic viscosity function was used in the Navier-Stokes equation for the mathematical simulations. Furthermore, as highlighted in a study of mussels by Riisgård and Larsen \cite{riisgardViscositySeawaterControls2007}, higher temperatures lead to low viscosity, which in term lead to higher cilia beating frequency. 

\subsection{Synchronisation}


%-------------------------------------------------------------
% END OF DOCUMENT
%-------------------------------------------------------------
\medskip

\printglossary

\printbibliography
\end{document}